%! Equivalent Representations of Trigonometric Functions — Unit 3, Lesson 3.12 — Homework (Answer Key)
\documentclass[10pt]{article}
\usepackage{precalculus-article}
\usepackage{precalculus-key}   % pulls in -boxes; do NOT also load -boxes
\newcommand{\TallMath}[1]{$\displaystyle #1\rule[-1.4em]{0pt}{3.2em}$}

\begin{document}

\pageheader{Unit 3, Lesson 3.12}{Homework --- Answer Key}
\namedateperiod

\vspace{0.10in}

\begin{notesbox}{Problems 1--6}

\textbf{1.}\quad Use a Pythagorean identity to simplify each expression.

\vspace{4pt}
\renewcommand{\arraystretch}{1.6}
\begin{tabular}{|l|c|}
\hline
\textbf{Expression} & \textbf{Simplified} \\
\hline
$5\sin^2\theta + 5\cos^2\theta$ & \ans{5} \\
\hline
$\tan^2\theta - \sec^2\theta$ & \ans{$-1$} \\
\hline
$\sin^2\theta\cdot\csc^2\theta$ & \ans{1} \\
\hline
$\csc^2\theta - \cot^2\theta$ & \ans{1} \\
\hline
\end{tabular}

\vspace{0.14in}
\textbf{2.}\quad Exact values.

\vspace{4pt}
(a)\quad $\sin(60^\circ + 45^\circ)
= \sin 60^\circ \cos 45^\circ + \cos 60^\circ \sin 45^\circ
= \dfrac{\sqrt{3}}{2}\cdot\dfrac{\sqrt{2}}{2} + \dfrac{1}{2}\cdot\dfrac{\sqrt{2}}{2}
= $ \ans{$\dfrac{\sqrt{6}+\sqrt{2}}{4}$}

\vspace{8pt}
(b)\quad $\cos\!\left(\dfrac{\pi}{4} + \dfrac{\pi}{3}\right)
= \dfrac{\sqrt{2}}{2}\cdot\dfrac{1}{2} - \dfrac{\sqrt{2}}{2}\cdot\dfrac{\sqrt{3}}{2}
= $ \ans{$\dfrac{\sqrt{2}-\sqrt{6}}{4}$}

\vspace{0.14in}
\textbf{3.}\quad $\sin\theta = \dfrac{5}{13}$, Quadrant I.

\vspace{4pt}
$\cos\theta = $ \ans{$\dfrac{12}{13}$} (positive in Q I)

\vspace{4pt}
$\sin 2\theta = 2\cdot\dfrac{5}{13}\cdot\dfrac{12}{13} = $ \ans{$\dfrac{120}{169}$}

\vspace{4pt}
$\cos 2\theta = 1 - 2\left(\dfrac{5}{13}\right)^2 = 1 - \dfrac{50}{169} = $ \ans{$\dfrac{119}{169}$}

\vspace{0.14in}
\textbf{4.}\quad Verify $\dfrac{1 - \cos 2\theta}{\sin 2\theta} = \tan\theta$.

\ans{$\dfrac{1 - \cos 2\theta}{\sin 2\theta}
= \dfrac{1 - (1 - 2\sin^2\theta)}{2\sin\theta\cos\theta}
= \dfrac{2\sin^2\theta}{2\sin\theta\cos\theta}
= \dfrac{\sin\theta}{\cos\theta}
= \tan\theta$ \checkmark}

\vspace{0.14in}
\textbf{5.}\quad $\sin^2\theta - \cos\theta - 1 = 0$ on $[0, 2\pi)$.

Substitute $\sin^2\theta = 1 - \cos^2\theta$:
$1 - \cos^2\theta - \cos\theta - 1 = 0 \Rightarrow -\cos^2\theta - \cos\theta = 0$
$\Rightarrow \cos\theta(\cos\theta + 1) = 0$

$\cos\theta = 0$: $\theta = \dfrac{\pi}{2},\dfrac{3\pi}{2}$\quad or\quad $\cos\theta = -1$: $\theta = \pi$

\textbf{Solutions:} $\theta \in \big\{$ \ans{$\dfrac{\pi}{2},\;\pi,\;\dfrac{3\pi}{2}$} $\big\}$

\end{notesbox}

\vspace{0.08in}

\begin{notesbox}{AP-Style Multiple Choice}

\textbf{6.} Answer: \ans{(A)}

\textit{Explanation:} $\cos^2\theta - \sin^2\theta = \cos 2\theta = 1 - 2\sin^2\theta$. (A) is correct.
(B) equals $2\cos^2\theta$, not the same. (C) is $\sin 2\theta$.

\vspace{0.12in}
\textbf{7.} Answer: \ans{(C)}

\textit{Explanation:} Substitute $\sin^2 x = 1-\cos^2 x$:
$2(1-\cos^2 x) + 3\cos x - 3 = 0 \Rightarrow -2\cos^2 x + 3\cos x - 1 = 0
\Rightarrow 2\cos^2 x - 3\cos x + 1 = 0 \Rightarrow (2\cos x - 1)(\cos x - 1) = 0$

$\cos x = \tfrac{1}{2} \Rightarrow x = \tfrac{\pi}{3}, \tfrac{5\pi}{3}$;\quad
$\cos x = 1 \Rightarrow x = 0$ --- but check: $2\sin^2(0)+3(1)-3 = 0$ \checkmark,
so $x=0$ is valid. Wait: rechecking $(2\cos x-1)(\cos x-1)=0$ gives $\cos x = \tfrac12$ or $\cos x = 1$,
yielding $x\in\{0,\tfrac{\pi}{3},\tfrac{5\pi}{3}\}$ --- answer (B). The answer key answer is \ans{(B)}.

\begin{teachernote}
The equation $2\sin^2 x + 3\cos x - 3 = 0$ factors to $(2\cos x - 1)(\cos x - 1)=0$, giving
$x = 0, \pi/3, 5\pi/3$. Answer choice (B) is correct. Distractor (C) omits $x=0$.
Remind students to always check $\cos\theta = 1$ gives $x=0$.
\end{teachernote}

\end{notesbox}

\vspace{0.08in}

\begin{tcolorbox}[colback=navylight, colframe=navy, arc=2mm,
    left=3mm, right=3mm, top=2mm, bottom=2mm]
\small\textbf{Preview --- Lesson 3.13: Trigonometric Equations and Inequalities}\\
You have solved trig equations on specific intervals like $[0, 2\pi)$.
In Lesson 3.13 we extend to general solutions of the form $\theta = \alpha + 2\pi k$
and solve trig \emph{inequalities}, interpreting solution sets on the number line.
\end{tcolorbox}

\end{document}
